\begin{graphillustration}{Пример: два треугольника и мост — три блока. Справа дерево блоков; двойным контуром отмечены точки сочленения.}{fig:05-block-tree}
\begin{tikzpicture}[font=\small]
\node[gi selected] (a) at (0,0) {$a$};
\node[gi vertex] (p) at (-1.2,1) {$p$};
\node[gi vertex] (q) at (-1.2,-1) {$q$};
\node[gi selected] (b) at (1.7,0) {$b$};
\node[gi vertex] (r) at (2.9,1) {$r$};
\node[gi vertex] (s) at (2.9,-1) {$s$};
\draw[gi edge] (a) -- (p);
\draw[gi edge] (p) -- (q);
\draw[gi edge] (q) -- (a);
\draw[gi edge] (a) -- (b);
\draw[gi edge] (b) -- (r);
\draw[gi edge] (r) -- (s);
\draw[gi edge] (s) -- (b);
\node[gi note] at (-0.9,1.65) {$B_1$};
\node[gi note] at (0.85,0.6) {$B_2$};
\node[gi note] at (2.7,1.65) {$B_3$};
\node[gi note] at (4.4,0) {$\longrightarrow$};
\begin{scope}[shift={(5.7,0)}]
\node[gi vertex] (B1) at (0,0) {$B_1$};
\node[gi selected] (a) at (1.4,0) {$a$};
\node[gi vertex] (B2) at (2.8,0) {$B_2$};
\node[gi selected] (b) at (4.2,0) {$b$};
\node[gi vertex] (B3) at (5.6,0) {$B_3$};
\draw[gi edge] (B1) -- (a);
\draw[gi edge] (a) -- (B2);
\draw[gi edge] (B2) -- (b);
\draw[gi edge] (b) -- (B3);
\end{scope}
\end{tikzpicture}
\end{graphillustration}
